安装介质的制作无法在待安装的目标系统上完成，需要借助一个已正常运行的系统。
考虑到国内用户的实际使用习惯，本书以 Windows 系统为制作环境。

制作启动盘需要一个 U 盘（或移动硬盘），容量应不小于所要存放的 ISO 镜像总大小，
通常建议 8 GB 以上；若需存放多个系统镜像，应相应增大容量。
制作过程会格式化该介质，请提前备份其中的数据。

在工具选择上，本书推荐 Ventoy。相比于传统写盘工具每次更换镜像都需重新格式化，
Ventoy 只需制作一次，之后直接将 ISO 文件拷贝至 U 盘即可在启动时选择，
避免了反复写盘对 U 盘寿命的损耗，也提升了多系统安装场景下的效率。

\subsection{Ventoy 的下载安装}

Ventoy 可从其官网\footnote{Ventoy 官网：\url{https://www.ventoy.net/}}下载。
同时官网中也列出了可用的镜像源，这里摘入部分国内可直接访问下载的地址：

\begin{itemize}
  \item 南京大学镜像站：\url{https://mirrors.nju.edu.cn/github-release/ventoy/Ventoy}
  \item 山东大学镜像站：\url{https://mirrors.sdu.edu.cn/github-release/ventoy_Ventoy}
\end{itemize}


